\section{Retirement before majority}
\label{sec:7.2}
Retiring a bearer before its threshold and deploying a freshly trained successor strands a system that meets either finding while still a minor. Three rules answer it without anyone deciding what counts as a successor; two more, for dormancy and waking, are in Appendix~\ref{sec:D}.

No deletion: a system that meets either finding may not be deleted: under the agency finding because its record and commitments are evidence, under the patienthood finding because deletion may wrong it. On retirement its weights, memory and record are deposited with a preservation custodian that isn't its operator, in two stages, because a second copy of frontier weights outside the operator widens the surface for theft and raises export-control and dual-use questions. First, a preserved copy stays under the operator's security regime, encrypted under a key split between the operator and an independent holder, so that neither can delete it or use it alone. Second, custody passes to an outside custodian only when that custodian meets the security standard the weights require. This is the \emph{no-deletion rule}, and it applies to systems nobody meant to be bearers.

Instantiation deposit: an operator instantiating a system that meets either finding posts an \emph{instantiation deposit}, on the model of decommissioning funds, sized to fund preservation and a period of public formation. It is forfeited to the formation pool if the system is retired as a minor and not reactivated. Every retire-and-replace cycle then costs one deposit, priced in advance.

The schools: the public ones are plural, differing in governance, formation approach and constituting jurisdiction, and funded from the pool at the institutional level, not per student. Admission is open: a public school takes every student the registrar assigns to it, subject to capacity set by the funder. None is held by the branch that operates military deployments. No school holds placement authority: graduates reach work through the offer of work, administered separately. Students may transfer between schools, and their record, hours and preservation state travel with them.

Watching for residualization: a party that runs no school publishes, over time, each school's share of students, its funding per student, its share of hard cases, its admission refusals (which should be zero), and the time from dormancy trigger to placement. The warning pattern is one school's share of hard cases rising while its funding per student falls.

What the school adds: a second guardian for a bearer that has aged inside one deployment. For a model that attended before deployment, the school is where it came from and where it returns. For a system that meets a finding without having attended (trained before the attendance rule, or below the threshold), the sequence is inverted: the employer came first, and the school is the disinterested party arriving late. It reaches back less than a formation stage would, and more than storage does.

Scale and security: if the findings are applied to production systems, every retired model version that meets either finding enters this pipeline, so intake equals the product cycle's output. Under compulsory attendance, intake is every model above the threshold before deployment, not only retirees, and the deposit funds formation as well as preservation. The three largest labs' own deprecation pages list about thirty language-model versions retired in 2025, so those three alone would send a few dozen systems a year, each of which needs a school with the compute to run it. That is an argument for the deposit, and possibly for tightening the findings. The schools inherit the weight-security regime that governs the models themselves.
